% ============================================================================================================
% BEFORE THIS DRAFT IS SHARED, CHECK (working list, not part of the paper; drafted 2 November 2026):
% (1) Every number against results/t91_carroll_reader_stage1.csv and _heldout.csv, PREREGISTRATION.md T91 with
%     Amendment 1 and its held-out reading, and ASSUMPTIONS.md O174. Table I and the modal counts of Table II were
%     regenerated from the two CSVs when this draft was written; recheck once.
% (2) Equation numbers of [CCM16] and [CC17] are those of the arXiv PDFs (1606.08444v3, 1712.02803v1) as recorded
%     in docs/reading/notes/2026-11-02_carroll_space_from_entanglement.md; check them against the journal versions.
%     The reading rule epsilon_D is unnumbered in the arXiv text (it sits between Eqs. (25) and (26)); the
%     docstring of scripts/carroll_reader.py calls it "Eq. 26", which is the stress function. Fix the docstring.
% (3) The [DT17] numbers of Sec. III.C come from dt17_calibration_check.py in this folder (log beside it), not from
%     the record's code. The reading note on Trugenberger's papers gives 4.0e-3 bits for pairs at distance 2 within
%     a copy; the exact values are 1.78e-3 (one shared neighbor) and 6.26e-3 (two), whose mean is 4.0e-3. The
%     result that the background-kept information graph has diameter 8.63 is new. Enter both in the record (an O
%     entry) before sharing.
% (3b) Sec. IV and V report a check made after the held-out reading (window_wrap_check.py, log beside it): the
%     registered R = 10 window on the 32 x 32 sheet wraps round the torus (pairs 20 apart > 16); on 64 x 64 it reads
%     3 (0.362, 0.081), as Amendment 1 predicted. Not in the record; PREREGISTRATION T91 and O174 should get a dated
%     note (the verdict as scored does not change). The reading note's section 7.2 rule "smaller than half the torus"
%     was not met at R = 10.
% (4) REFERENCES.bib: [DT17] still says "abstract only" while the reading note of 2 November reads it in full (the
%     note's replacement line is owed); its volume and page are to verify. [T17]'s note says not read; the same
%     reading note reads it in full. [CC17] and [CS20] journal references are from the arXiv listing; verify.
% (5) The owner's paragraph in Sec. I (her hypothesis and her positions of VISION Update 152): her wording to be
%     confirmed by her.
% (6) The registered threshold (0.1 of the largest eigenvalue) is ours; the flat sheet's fourth eigenvalue sits at
%     0.101 of the largest at R = 4 and 10, on the threshold. Say so wherever a count is quoted.
% (7) Figures this draft will need (drawn by script from results/t91_*.csv, none invented): (a) eigenvalue shares
%     against window width 2R + 1, log-log, for the flat sheet (third and fourth) and the tube (curled pair), with
%     the R = 4, 6 fit and the held-out R = 10 point and its range; (b) the distribution of D_MDS over centers by R
%     for the opened sheets and the melts.
% (8) A fresh red-team pass (CLAUDE.md rule 12) and a check that no sentence says "not studied" without a note
%     (rule 15). No physicist has read this.
% ============================================================================================================
\documentclass[aps,pre,10pt,twocolumn,superscriptaddress,nofootinbib]{revtex4-2}
\usepackage{graphicx}
\usepackage{amsmath}
\usepackage{booktabs}
\usepackage[hidelinks]{hyperref}

\begin{document}

\title{Does entanglement geometry see the same space? The mutual-information reader of Cao, Carroll and
Michalakis applied to the wirings of a graph model of emergent geometry}

\author{Emily Smith}
\email{emilysmithcreate@gmail.com}
\affiliation{Independent researcher, Charlottesville, Virginia, USA}

\date{DRAFT, 2 November 2026}

\begin{abstract}
Two published routes build space from a network. In combinatorial quantum gravity (CQG) the network is a classical
wiring weighted by a discrete curvature; in the program of Cao, Carroll and Michalakis it is a quantum state on a
factored Hilbert space, and distances come from the mutual information between factors. We ask what the second
route's reader reports when it is pointed at the first route's wirings. With one factor per vertex and the same
mutual information on every link, their distance ansatz with $\Phi=-\ln$ and its shortest-path completion give the
graph distance up to a constant (exact), so their dimension estimator becomes classical multidimensional scaling
(MDS) of graph-distance windows. Whole tori cannot be read: each wrapped direction takes a pair of large
eigenvalues (exact). On balls of graph radius $R=2$ to 10, with a count fixed before the computation (eigenvalues at
least 0.1 of the largest) and $R=10$ held out: the 4-cube, which the model's local count calls zero-dimensional,
reads four; a flat square lattice reads three to four, not two, because its graph distance is a taxicab distance; a
direction curled to four sites is counted at small windows and its share falls as $(2R+1)^{-1.24}$, a law fitted at
$R=4,6$ and confirmed at $R=10$ ($0.112$ against a predicted $0.076$ to $0.141$); and the narrow tori that the
model's conversion produces read one at $R=10$. Three of five registered predictions failed at one or more radii.
No quantum state of the model is computed; one published state built from a wiring is used as a calibration.
\end{abstract}

\maketitle

\section{Introduction}

There are at least two published ways to get space out of a network. In combinatorial quantum gravity (CQG)
\cite{T17,KTB19,T25} the network is a classical graph and its energy is the total Ollivier curvature of the wiring;
at low coupling the favored graphs are lattice-like and a geometry is read from the wiring itself. In the program
of Cao, Carroll and Michalakis \cite{CCM16,CC17} the network is derived: a quantum state on a Hilbert space factored
into many pieces defines a weighted graph whose weights are the mutual informations between pieces, a distance is
assigned to each weight, and a dimension is read by embedding the distances in Euclidean space. Ref.~\cite{CCM16}
says plainly that it assumes no Hamiltonian and that the factorization is given; Ref.~\cite{CS20} asks how a
factorization could be chosen.

The two routes meet at one point: both end with a weighted graph and ask what space it describes. This paper asks a
narrow question at that point. \emph{When the second route's reader is given the first route's wirings, does it see
the same space?} In particular, does it agree with the local count of dimension that a square-counting model uses,
on a flat sheet, on a sheet curled in one direction, and on a fully curled object?

The wirings come from a family of energies next to CQG, studied in a companion paper \cite{paper1} (Sec.~\ref{sec:toy}).
The reader is the one of Ref.~\cite{CCM16} without change, except for one choice that applying it to a classical
wiring requires (Sec.~\ref{sec:reader}). The computation was registered before it was run, with its predictions
and one size held out \cite{prereg}; Sec.~\ref{sec:results} reports it as registered, including the predictions that
failed.

\emph{Why the author asks (her hypothesis, stated once).} The author holds that our space is a stable phase of a
deeper, specific arrangement that is curled, and that space arose when it opened, releasing a bounded amount of
energy. On the parameter used below she holds that $\lambda$ is 1 or below everywhere when the network is taken as a
whole, the curled object seen one part at a time being the exception; that the two sides of the bipartite graph are
two distinct kinds of point, every point within a kind the same; and that points belong to the two-dimensional
description only, the foundational degrees of freedom being twist and curl. Whether a curled direction reads as a
dimension is the question her picture of opening rests on, and that is why the reader's answer matters to her. None
of this is assumed below.

\emph{Kinds of statement.} Each result is marked as exact (arithmetic or complete enumeration), measured (a computed
reading registered before it was made, with its verdict), or ours and unverified (an interpretation). Nothing in
this paper is about the real universe.

\section{The model, briefly}
\label{sec:toy}

A state is a four-regular bipartite graph on $N$ vertices obeying the hard-core rule (no two vertices share more than
two neighbors) \cite{KTB19,T25}. With $S$ the number of squares (four-cycles), $S_e$ the number on edge $e$, and
$X=\sum_e(S_e-2)_+$, the energy is
\begin{equation}
H = 16\,(N - S) + 4\lambda X .
\label{eq:H}
\end{equation}
At $\lambda=1$ this is the CQG curvature action, which Ref.~\cite{T25} (its Eq.~(21)) writes as
\[
H=-4\sum_v\sum_{w\sim v}\kappa(v,w),
\]
with $\kappa$ the Ollivier curvature of the link. Only $\lambda=1$ is CQG; every other value is our family around it
\cite{paper1}.

\emph{The local count (exact).} A vertex has six pairs of incident edges; under the hard-core rule a pair closes at
most one square. Let $d(v)$ be the number of pairs that close none. On the flat lattice torus $d=2$; on the
$4\times L$ torus, one direction curled to a single square, $d=1$; on the $4\times4$ torus, which is the 4-cube,
$d=0$ \cite{paper1}. This is a count of square-free directions at one vertex, not a certified geometric dimension.
For $\lambda>1$ the curled torus lies $4(\lambda-1)$ per vertex above the flat one, and after activation it converts
with that release; the converted states measured so far are narrow tori whose shortest loop around the short
direction is six edges \cite{paper1}.

\section{The reader}
\label{sec:reader}

\subsection{Definitions, as published}

Ref.~\cite{CCM16} factors the Hilbert space into finite-dimensional pieces, $\mathcal H=\bigotimes_p\mathcal H_p$,
and uses the mutual information of two pieces,
\[
I(A_p:A_q)=S(A_p)+S(A_q)-S(A_pA_q).
\]
A state is redundancy-constrained when the entropy of every union $B$ of pieces is a cut sum, which
Ref.~\cite{CCM16} (its Eq.~(9)) writes as
\[
S(B)=\tfrac12\sum_{p\in B,\;q\notin B}I(A_p:A_q).
\]
Each pair with nonzero mutual information gets an edge of length, in their Eq.~(13),
\begin{equation*}
w(p,q)=\ell_{\rm RC}\,\Phi\!\left(\frac{I(A_p:A_q)}{I_0}\right),
\end{equation*}
with $\Phi(1)=0$, $\Phi(x)\to\infty$ as $x\to0$, $\Phi$ decreasing, and $I_0$ the mutual information of a maximally
entangled pair; their examples use $\Phi=-\ln$. No edge is drawn where $I=0$. The distance is the shortest path,
their Eq.~(14),
\begin{equation*}
d(p,q)=\min_{P}\sum_{n}w(p_n,p_{n+1}).
\end{equation*}

The dimension is read by classical multidimensional scaling (their Sec.~III.C). For $n$ points with distances
$d(p,q)$ the centered matrix of their Eq.~(23) is
\begin{equation*}
\begin{aligned}
B_{pq}=-\tfrac12\Bigl[\,&d(p,q)^2-\tfrac1n\sum_l d(p,l)^2\\
&-\tfrac1n\sum_l d(l,q)^2+\tfrac1{n^2}\sum_{l,m}d(l,m)^2\Bigr],
\end{aligned}
\end{equation*}
diagonalized as $B=V\Lambda V^{\mathsf T}$ with eigenvalues $\lambda_1\ge\lambda_2\ge\dots$ (their Eq.~(24)); the
embedding in $D$ dimensions uses the top $D$ eigenvectors scaled by $\sqrt{\lambda_i}$ (their Eq.~(25)). If the
distances are exactly Euclidean in $D$ dimensions, $B$ has exactly $D$ positive eigenvalues and the rest zero. If not,
they read the dimension from the spectrum with the rule, stated after their Eq.~(25),
\begin{equation*}
\varepsilon_D=1-\frac{\sum_{i=1}^{D}|\lambda_i|}{\sum_{i=1}^{n}|\lambda_i|},
\end{equation*}
an embedding in $D$ dimensions being ``approximately valid'' when $\varepsilon_D$ is small. Their own two examples
give a scale for ``small'': a spin-1 Heisenberg chain of 100 sites, $\varepsilon_1=0.0167$, and a coarse-grained
toric-code patch, $\varepsilon_2=0.41$, called evidently two-dimensional \cite{CCM16}.

\subsection{The one choice: a factor per vertex, the same mutual information on every link}

To apply the reader to a classical wiring two things must be chosen.

\emph{A factor per vertex.} Their factors are the vertices of the information graph and the mutual informations are
its edges. The model's vertices are what its links join, so the reading that changes least is one factor per vertex
and a nonzero mutual information exactly on the links. A factor per link would read the line graph instead and give
different distances; it is not done here. With the author's position that points belong only to a two-dimensional
description (Sec.~I), this is a reading of that description, not of the degrees of freedom she takes as
foundational.

\emph{A mutual information per link.} The model says only whether two vertices are linked; all links are alike. So
take $I(A_p:A_q)/I_0=x_0$ on every link, for one constant $0<x_0<1$, and zero otherwise.

\emph{Exact.} Then every edge has the same length $w=\ell_{\rm RC}(-\ln x_0)$, Eq.~(14) gives $d(p,q)=w\,\delta(p,q)$
with $\delta$ the graph distance, and $B$ is $w^2$ times the matrix built from $\delta$. Its eigenvalues scale by $w^2$,
so their ratios, $\varepsilon_D$ and any count of eigenvalues relative to the largest do not depend on $x_0$,
$\ell_{\rm RC}$ or $I_0$. ($x_0=1$ is excluded: it gives $w=0$ and every point at one place.) The reader of a wiring
is therefore classical MDS of graph distances. Their toric-code example is close to this case: there the mutual
information on each link of a uniform coarse-graining is a boundary length up to a constant \cite{CCM16}.

A choice that weights each link by its square count would put the model's energy into the lengths; that is a second,
different reading and is not made here.

\subsection{One explicit state whose mutual information follows a wiring}

The constant above is put in by hand. To show that a quantum state with such mutual information exists, and to see
what else such a state brings, we use the one state in the CQG literature built from a wiring. Ref.~\cite{DT17} (its
Eq.~(4)) turns the adjacency matrix $w_{ij}$ of a graph into a pure state on two copies of an $N$-level register,
\begin{equation*}
|w\rangle=\frac{1}{\lVert w\rVert_F}\sum_{i,j}w_{ij}\,|i\rangle|j\rangle ,
\end{equation*}
which it reads equivalently as two registers of $N$ qubits, each with exactly one excitation. Ref.~\cite{DT17}
cites Ref.~\cite{CCM16} for reading space as an entanglement pattern, but it computes one spectral number per
graph, not a geometry.

\emph{Exact (our computation, 8 $\times$ 8 flat torus, one factor per qubit, $2N=128$ factors).} Between qubit $i$ of
one copy and qubit $j$ of the other the mutual information is $1.17\times10^{-2}$ bits when $i$ and $j$ are linked and
$3.58\times10^{-4}$ bits otherwise. Within one copy it is $1.78\times10^{-3}$ bits for vertices sharing one neighbor,
$6.26\times10^{-3}$ for vertices sharing two, and $3.58\times10^{-4}$ for all other pairs. So the mutual information
follows the wiring across the copies, and its distance-two graph within each.

\emph{Exact (same state).} The uniform background is nonzero, so Eq.~(13) draws an edge between every pair. With
$\Phi=-\ln$ and $I_0=2$ bits, a link has length $5.14\,\ell_{\rm RC}$ and a background edge $8.63\,\ell_{\rm RC}$,
less than two links. The shortest path therefore never uses more than one background edge, every distance is at most
$8.63\,\ell_{\rm RC}$, and MDS of all 128 factors gives one eigenvalue and then five equal ones at 0.76 of it: no
low-dimensional geometry. With the background removed the graph splits into two pieces of 64 factors, each with
the torus's wiring (bipartite, so the two copies pair one side of one copy with the other side of the other) plus
the distance-two edges, and each reads like a whole torus (four equal eigenvalues; Sec.~\ref{sec:windows}).

\emph{Ours, unverified.} The state shows that the mutual information of a quantum state can follow a wiring, and it
shows that it does so only up to a background that, read literally, erases the geometry. It is built from the wiring,
so it calibrates the reader and tests nothing about how a geometry could arise. Whether it is redundancy-constrained
was not checked.

\section{Why whole tori cannot be read, and windows}
\label{sec:windows}

\emph{Exact.} On a torus the distance matrix is unchanged by translations, so $B$'s eigenvectors are Fourier modes,
which come in equal pairs, a cosine and a sine. Each wrapped direction's longest mode takes one such pair. The
$8\times8$ torus gives four equal eigenvalues, 218.5, then four at 37.5; the $4\times16$ curled torus gives 525.5,
525.5 (the long direction), then 160, 160 (the curled one, 0.30 of the first and growing with the length of the
torus), then 64.8, 64.8; the 4-cube gives exactly four eigenvalues of 16, the others 0 or $-4$. Even a single square,
the curled direction's cross-section, is not Euclidean under graph distance (opposite corners at 2, not $\sqrt2$),
and its spectrum is $2,2,0,-1$. So on whole tori the sheet, the curled torus and the 4-cube all show four large
eigenvalues, distinguished by ratios that depend on the torus's size rather than its wiring. Ref.~\cite{CCM16}
reads open patches only (a chain and a toric-code patch).

\emph{Windows.} The reader is applied to a ball of graph radius $R$ around a center vertex, using the whole graph's
shortest-path distances between the ball's vertices. On the flat $32\times32$ torus no path inside the ball wraps
for $R\le8$; at $R=10$ the ball holds pairs 20 apart along one axis, more than half of 32, and some of their shortest
paths wrap (this was not noticed when the computation was registered; Sec.~\ref{sec:results}). The curled direction
of the tube wraps inside any ball with $R\ge2$, and so does the short direction of the converted narrow tori (six
around) for $R\ge3$; there the wrap cannot be cut out and is part of what is read.

\emph{The registered count.} Ref.~\cite{CCM16} leaves ``small'' open. Before the computation we fixed
\[
D_{\rm MDS}=\#\{\,i:\lambda_i\ge0.1\,\lambda_1\,\}
\]
and record $\varepsilon_1$ to $\varepsilon_4$ beside it. The threshold is ours.

\section{Results}
\label{sec:results}

\subsection{What was registered}

The states \cite{prereg}: built (exact, no sampling) the flat $32\times32$ torus, the $4\times64$ curled torus and the
4-cube, each read from one center (every center is equivalent); measured states saved by earlier runs, namely four
converted narrow tori of 1,024 vertices from the slowest-cooled runs at $\lambda=1.25$ (every eighth vertex a center,
512 windows per radius) and four warm random graphs of the published model ($\lambda=1$, $N=196$, $g=3.3$; every vertex
a center, 784 windows). Radii $R=2,3,4,6$; $R=10$ held out, to be computed only after a written prediction from the
smaller radii. The predictions were ours (the author gave none):
\begin{itemize}
\item P0, the 4-cube reads $D_{\rm MDS}=4$ for $R\ge4$, against $d=0$;
\item P1, the flat sheet reads 2 at $R=3,4,6$;
\item P2, the curled torus reads no more than the sheet at $R=4,6$, and its curled pair's share
$\lambda_2/\lambda_1$ is near $6/(2R+1)^2$ within a factor of two (0.074 at $R=4$, 0.036 at $R=6$), so that at
$R\ge4$ the curled direction is not counted;
\item P3, the converted tori's most common count equals the flat sheet's at every $R$;
\item P4, the random graphs read at least 4 at $R=3$ at most centers.
\end{itemize}

\begin{table}[t]
\caption{Built states, exact. Eigenvalues of $B$ relative to the largest, the registered count and the reading rule,
by window radius. The toy's local count is $d=2$ (sheet), 1 (curled torus), 0 (4-cube) at every vertex. For
$R\ge4$ the 4-cube window is the whole cube.}
\label{tab:built}
\begin{ruledtabular}
\begin{tabular}{lrrccccr}
State & $R$ & $n$ & $\lambda_2/\lambda_1$ & $\lambda_3/\lambda_1$ & $\lambda_4/\lambda_1$ & $\varepsilon_2$ &
$D_{\rm MDS}$\\
\hline
Sheet & 2 & 13 & 1.000 & 0.521 & 0.164 & 0.38 & 4\\
 & 3 & 25 & 1.000 & 0.428 & 0.123 & 0.41 & 4\\
 & 4 & 41 & 1.000 & 0.396 & 0.101 & 0.42 & 4\\
 & 6 & 85 & 1.000 & 0.373 & 0.087 & 0.43 & 3\\
 & 10 & 221 & 1.000 & 0.408 & 0.101 & 0.47 & 4\\
Curled & 2 & 12 & 0.607 & 0.594 & 0.093 & 0.44 & 3\\
 & 3 & 20 & 0.426 & 0.395 & 0.062 & 0.42 & 3\\
 & 4 & 28 & 0.309 & 0.291 & 0.053 & 0.39 & 3\\
 & 6 & 44 & 0.196 & 0.189 & 0.041 & 0.33 & 3\\
 & 10 & 76 & 0.112 & 0.110 & 0.027 & 0.25 & 3\\
4-cube & 2 & 11 & 1.000 & 1.000 & 0.252 & 0.54 & 4\\
 & 3 & 15 & 1.000 & 1.000 & 0.702 & 0.61 & 4\\
 & 4+ & 16 & 1.000 & 1.000 & 1.000 & 0.64 & 4\\
\end{tabular}
\end{ruledtabular}
\end{table}

\subsection{The smaller radii (measured)}

The built states are in Table~\ref{tab:built}.

\emph{P0 holds.} The 4-cube reads four at every radius; the model's count reads zero.

\emph{P1 fails.} The flat sheet reads 4, 4, 4, 3 at $R=2,3,4,6$. Its third eigenvalue is 0.37 to 0.52 of the
largest, and its fourth sits near the threshold (0.101 at $R=4$). Its $\varepsilon_2$ is 0.38 to 0.43, the same
level as the toric-code patch that Ref.~\cite{CCM16} calls two-dimensional; the registered count, not the spectrum's
overall spread, is what calls it three or four.

\emph{P2 holds on the ranking and fails on the size.} The curled torus reads 3 against the sheet's 4 at $R=4$ and 3
against 3 at $R=6$. Its curled pair is 0.309 and 0.196 of the largest at $R=4$ and 6, four and five times the
predicted 0.074 and 0.036, so at these radii the curled direction is counted.

\emph{P3 fails at $R=2$ only.} The converted tori read 5 at $R=2$ (508 of 512 windows) against the sheet's 4, and the
sheet's count at $R=3,4,6$ (4, 4, 3 in at least 507 of 512). Their windows are not isotropic ($\lambda_2/\lambda_1$
averages 0.50 at $R=2$, against 1.00 for the square sheet): the narrow torus shows.

\emph{P4 holds.} The random graphs read at least 5 at 778 of 784 centers at $R=3$, and the count grows with $R$
(5 to 16 at $R=4$; 9 to 17 at $R=6$).

\subsection{The held-out radius}

\emph{Prediction, written before $R=10$ was computed} \cite{prereg}. Fitting power laws in the window's width
$2R+1$ to $R=4$ and 6 only: the curled pair's share at $R=10$ would be 0.108 (exponent 1.24), with a range of
$\pm30\%$, 0.076 to 0.141; the sheet's third eigenvalue 0.345 (0.241 to 0.448) and fourth 0.071. Predicted counts:
sheet 3, curled torus 3 (its 0.108 sits just above the threshold, so 1 was named as possible), 4-cube 4.

\emph{Measured.} The curled pair is 0.112 and 0.110 of the largest, inside the predicted range: the size claim holds
at the held-out radius. Extrapolating the fit (not computed), the share crosses 0.1 at windows about 22 sites wide.
The curled torus reads 3 and the 4-cube 4, as predicted. The sheet reads 4, not the predicted 3: its third eigenvalue,
0.408, is inside its range, and its fourth, 0.101, is above the predicted 0.071 and on the threshold. \emph{A check
made after this reading, not registered (exact):} the same window on a $64\times64$ torus, where no path inside it
wraps, gives 0.362 and 0.081 and reads 3, while the curled torus reads the same at $4\times128$ as at $4\times64$. So
the sheet's miss at $R=10$ comes from paths wrapping round the $32\times32$ torus, a flaw of the registered design;
the registered verdict stands as scored.

\emph{Seen, not predicted for this radius.} The converted narrow tori read 1 at 506 of 512 centers at $R=10$
(Table~\ref{tab:saved}): their windows wrap the short direction, as the curled torus's windows wrap its curled one,
and at this scale the reader calls them one-dimensional. P3 therefore fails at $R=10$ as well.

\begin{table}[t]
\caption{Saved states, measured. Most common $D_{\rm MDS}$ over the centers (with its share), by radius. Converted
tori: four 1,024-vertex products of the slowest-cooled runs at $\lambda=1.25$, 512 windows per radius; the model's
count is $d=2$ at 509 of 512 centers. Random graphs: four warm states of the published model, $N=196$, 784 windows;
at $R=10$ each window is the whole graph.}
\label{tab:saved}
\begin{ruledtabular}
\begin{tabular}{rcc}
$R$ & Converted tori & Random graphs (range)\\
\hline
2 & 5 (508/512) & 4 (439/784; 3 to 6)\\
3 & 4 (509/512) & 6 and 7 (260 each; 4 to 9)\\
4 & 4 (507/512) & 9 (176/784; 5 to 16)\\
6 & 3 (507/512) & 13 (200/784; 9 to 17)\\
10 & 1 (506/512) & 13 (392/784; 13 to 16)\\
\end{tabular}
\end{ruledtabular}
\end{table}

\emph{Verdict.} By the registered words, MIXED: P0 and P4 hold; P1 fails; P2 holds on the ranking, fails on the size
at $R=4,6$, and its held-out size claim holds at $R=10$; P3 fails at $R=2$ and $R=10$.

\section{Discussion}
\label{sec:discussion}

\emph{The taxicab artifact.} \emph{Exact:} on a square lattice the graph distance between vertices separated by
$(\Delta x,\Delta y)$ is $|\Delta x|+|\Delta y|$, so
\[
\delta^2=\Delta x^2+\Delta y^2+2|\Delta x||\Delta y| .
\]
The first two terms give the two eigenvalues of a plane; the last is not of the form that MDS can absorb into two
coordinates, and it produces the extra eigenvalues. \emph{Measured:} the third eigenvalue stays between 0.36 and 0.43 of
the largest from $R=3$ to 10 (0.362 at $R=10$ on the larger torus) and does not fade, and at $R=4$ and 6 negative
eigenvalues carry about a fifth of $\sum|\lambda_i|$. So
under a constant mutual information per link the reader does not call a square lattice two-dimensional at any
window tried, and this says more about the distance assigned to the lattice than about the model. \emph{Exact
arithmetic:} if instead every pair had $I/I_0=e^{-r_{pq}/\xi}$ with $r_{pq}$ the Euclidean distance between lattice
sites, then $w=(\ell_{\rm RC}/\xi)\,r_{pq}$, which already satisfies the triangle inequality, so Eq.~(14) returns it
unchanged and $B$ has exactly two positive eigenvalues on any patch of the plane. \emph{Ours, unverified:} a state
with that property needs mutual information between unlinked pairs that knows the Euclidean distance, which the
wiring alone does not supply; whether natural states on a lattice (gapped ground states, for which Ref.~\cite{CCM16}
suggests $\Phi=-\ln$) come close enough is the question.

\emph{A short curled direction depends on scale.} The curled direction's share falls as $(2R+1)^{-1.24}$ (measured,
confirmed at the held-out radius), more slowly than the $(2R+1)^{-2}$ of P2's estimate, which treated the four-site
loop as a fixed small circle beside a growing line. \emph{Ours, unverified:} in this reader a direction curled to
four sites counts at windows narrower than about twenty sites and stops counting beyond, which is the Kaluza-Klein
picture of a small direction that is invisible at large scale. The reader's dimension is then a function of scale,
like the spectral dimension reviewed in Ref.~\cite{Carlip17}, while the model's count $d$ is a one-step count that
calls the tube one-dimensional at every scale.

\emph{Fully curled reads four.} The 4-cube, which the model's count calls zero-dimensional, embeds with four equal
eigenvalues at every radius (exact). The two conventions part here by definition: the model counts the directions
along which no square closes, and the reader counts the independent directions of a distance embedding. A fully
curled object has many directions, all short.

\emph{The conversion's product reads one at large windows.} The converted tori are six sites around (Sec.~\ref{sec:toy}).
Their windows read 5, 4, 4 and 3 at $R=2$ to 6 and 1 at $R=10$, while the model's count reads 2 at nearly every
vertex. \emph{Ours, unverified:} by the reader's scale-dependent reading these products are one-dimensional at large
scale, with a short closed direction, which agrees with the companion paper's finding that the conversion does not
produce two large directions \cite{paper1}.

\section{Questions for the program}
\label{sec:questions}

\begin{enumerate}
\item How should the estimator treat a lattice whose graph distance is not Euclidean, as the square lattice's is not?
Is the intended reading the registered count, $\varepsilon_D$ against an example's value, or a change in $\Phi$ or in
the mutual information assigned to unlinked pairs? And does a direction curled to a few sites count as a dimension at
a given scale, or is the entropy-of-a-ball exponent of Ref.~\cite{CCM16} (its Eq.~(15)), which does not grow along
such a direction, the estimator the program intends when small closed directions are present?
\item Which factorization should be used for a classical wiring: one factor per vertex, one per link, or one chosen
by a dynamics as in Ref.~\cite{CS20}? And which weighted graphs are the information graphs of some
redundancy-constrained state, given that the one state built from a wiring here carries a background that, read
literally, removes the geometry?
\item Ref.~\cite{CC17} suggests that lattice models might be ``gravitized'' by emerging their geometry from a state's
entanglement rather than from the lattice. Is a classical wiring with a curvature energy, such as CQG, a candidate
for this, that is, is there a natural state for each wiring whose reader curvature (Ref.~\cite{CCM16}, its Eqs.~(36)
and (40)) tracks the wiring's Ollivier curvature?
\end{enumerate}

\section{What this paper does not show}

No quantum state of the model is computed. The mutual information used is a constant put on the links; the one
explicit state used is built from the wiring and serves only as a calibration. The readings depend on the registered
threshold, which is ours, and the flat sheet's count sits on it at two radii; at $R=10$ the registered sheet window
also wrapped round its torus. The converted tori come from runs at
$\lambda=1.25$, outside CQG; the random graphs are of the published model at $\lambda=1$; the built states are
geometry only. Agreement or disagreement between the reader and the model's count is a relation between two
definitions on the same graphs, not evidence about either route's account of space, and nothing here is about the
real universe.

\begin{acknowledgments}
The analysis, code and text were drafted with Claude (Anthropic), in conversation with the author, who directed and
checked the work. No physicist has reviewed it. The repository holds the computation, its registered predictions and
results: \url{https://github.com/EmilySmithCreate/RecreationalPhysics}.
\end{acknowledgments}

\begin{thebibliography}{99}
\bibitem{T17} C.~A. Trugenberger, Combinatorial quantum gravity: geometry from random bits, JHEP \textbf{09}, 045
(2017), arXiv:1610.05934.
\bibitem{KTB19} C.~Kelly, C.~A. Trugenberger, and F.~Biancalana, Self-assembly of geometric space from random graphs,
Class. Quantum Grav. \textbf{36}, 125012 (2019), arXiv:1901.09870.
\bibitem{T25} C.~A. Trugenberger, Networks as the fundamental constituents of the universe, J. Phys. Complex.
\textbf{6}, 042001 (2025), arXiv:2512.17676.
\bibitem{CCM16} C.~Cao, S.~M. Carroll, and S.~Michalakis, Space from Hilbert space: recovering geometry from bulk
entanglement, Phys. Rev. D \textbf{95}, 024031 (2017), arXiv:1606.08444.
\bibitem{CC17} C.~Cao and S.~M. Carroll, Bulk entanglement gravity without a boundary: towards finding Einstein's
equation in Hilbert space, Phys. Rev. D \textbf{97}, 086003 (2018), arXiv:1712.02803.
\bibitem{CS20} S.~M. Carroll and A.~Singh, Quantum mereology: factorizing Hilbert space into subsystems with
quasi-classical dynamics, Phys. Rev. A \textbf{103}, 022213 (2021), arXiv:2005.12938.
\bibitem{paper1} E.~Smith, The curled torus burps: activated escape and conversion in a graph model of emergent
geometry, technical manuscript (2026), \texttt{docs/papers/curled\_torus/} in the repository below.
\bibitem{prereg} E.~Smith, \texttt{PREREGISTRATION.md}, section T91 with Amendment 1, and \texttt{ASSUMPTIONS.md},
entry O174, in the repository below (2 November 2026); computation \texttt{scripts/carroll\_reader.py}, results
\texttt{results/t91\_carroll\_reader\_stage1.csv} and \texttt{\_heldout.csv}.
\bibitem{DT17} M.~C. Diamantini and C.~A. Trugenberger, Topological network entanglement as order parameter for the
emergence of geometry, New J. Phys. (2017), doi:10.1088/1367-2630/aa8f08, arXiv:1710.00221.
\bibitem{Carlip17} S.~Carlip, Dimension and dimensional reduction in quantum gravity, Class. Quantum Grav.
\textbf{34}, 193001 (2017), arXiv:1705.05417.
\end{thebibliography}

\end{document}
